% Options for packages loaded elsewhere
\PassOptionsToPackage{unicode}{hyperref}
\PassOptionsToPackage{hyphens}{url}
\PassOptionsToPackage{dvipsnames,svgnames,x11names}{xcolor}
%
\documentclass[
  11pt,
  a4paper,
]{ltjsarticle}
\usepackage{amsmath,amssymb}
\usepackage{iftex}
\ifPDFTeX
  \usepackage[T1]{fontenc}
  \usepackage[utf8]{inputenc}
  \usepackage{textcomp} % provide euro and other symbols
\else % if luatex or xetex
  \usepackage{unicode-math} % this also loads fontspec
  \defaultfontfeatures{Scale=MatchLowercase}
  \defaultfontfeatures[\rmfamily]{Ligatures=TeX,Scale=1}
\fi
\usepackage{lmodern}
\ifPDFTeX\else
  % xetex/luatex font selection
\fi
% Use upquote if available, for straight quotes in verbatim environments
\IfFileExists{upquote.sty}{\usepackage{upquote}}{}
\IfFileExists{microtype.sty}{% use microtype if available
  \usepackage[]{microtype}
  \UseMicrotypeSet[protrusion]{basicmath} % disable protrusion for tt fonts
}{}
\makeatletter
\@ifundefined{KOMAClassName}{% if non-KOMA class
  \IfFileExists{parskip.sty}{%
    \usepackage{parskip}
  }{% else
    \setlength{\parindent}{0pt}
    \setlength{\parskip}{6pt plus 2pt minus 1pt}}
}{% if KOMA class
  \KOMAoptions{parskip=half}}
\makeatother
\usepackage{xcolor}
\usepackage[margin=24mm]{geometry}
\setlength{\emergencystretch}{3em} % prevent overfull lines
\providecommand{\tightlist}{%
  \setlength{\itemsep}{0pt}\setlength{\parskip}{0pt}}
\setcounter{secnumdepth}{5}
\ifLuaTeX
  \usepackage{selnolig}  % disable illegal ligatures
\fi
\IfFileExists{bookmark.sty}{\usepackage{bookmark}}{\usepackage{hyperref}}
\IfFileExists{xurl.sty}{\usepackage{xurl}}{} % add URL line breaks if available
\urlstyle{same}
\hypersetup{
  colorlinks=true,
  linkcolor={blue},
  filecolor={Maroon},
  citecolor={Blue},
  urlcolor={blue},
  pdfcreator={LaTeX via pandoc}}

\author{}
\date{}

\begin{document}

\hypertarget{ux6f14ux7fd2-09-ux4f4eux30e9ux30f3ux30afux8fd1ux4f3cux3068ux30c7ux30fcux30bfux89e3ux6790}{%
\section{演習 09 ---
低ランク近似とデータ解析}\label{ux6f14ux7fd2-09-ux4f4eux30e9ux30f3ux30afux8fd1ux4f3cux3068ux30c7ux30fcux30bfux89e3ux6790}}

\hypertarget{rank-k-ux8fd1ux4f3c}{%
\subsection{1. rank-k 近似}\label{rank-k-ux8fd1ux4f3c}}

特異値が \texttt{(10,3,1,0.1)} の行列について rank-2 近似の Frobenius
誤差を求めてください。

\hypertarget{ux30a8ux30cdux30ebux30aeux30fcux4fddux6301ux7387}{%
\subsection{2.
エネルギー保持率}\label{ux30a8ux30cdux30ebux30aeux30fcux4fddux6301ux7387}}

同じ特異値について、rank-2 近似が
\texttt{\textbar{}\textbar{}A\textbar{}\textbar{}\_F\^{}2}
の何割を保持するか求めてください。

\hypertarget{pca}{%
\subsection{3. PCA}\label{pca}}

中心化データ \texttt{X\_c} に対して共分散行列が
\texttt{(1/(n-1))X\_c\^{}TX\_c} となることを確認し、SVD
から主成分方向を導いてください。

\hypertarget{pca-ux306eux4e8cux3064ux306eux5b9aux7fa9}{%
\subsection{4. PCA
の二つの定義}\label{pca-ux306eux4e8cux3064ux306eux5b9aux7fa9}}

「射影後の分散最大化」と「再構成誤差最小化」が同じ主成分を与える理由を説明してください。

\hypertarget{ux30ceux30a4ux30baux9664ux53bb}{%
\subsection{5. ノイズ除去}\label{ux30ceux30a4ux30baux9664ux53bb}}

truncated SVD で \texttt{k}
を小さくしすぎる場合と大きくしすぎる場合に、それぞれ何が起こるか説明してください。

\end{document}
